\subsubsection{Binomial Approximation}
\label{sec:approximation}

\begin{figure*}[!ht]
    \centering
    \resizebox*{1\textwidth}{!}{% This file was created with tikzplotlib v0.10.1.
\begin{tikzpicture}

\definecolor{darkgray176}{RGB}{176,176,176}
\definecolor{darkgreen}{RGB}{0,100,0}
\definecolor{darkorange}{RGB}{255,140,0}
\definecolor{darkslateblue10251153}{RGB}{102,51,153}
\definecolor{green}{RGB}{0,128,0}
\definecolor{orange}{RGB}{255,165,0}
\definecolor{purple}{RGB}{128,0,128}

\begin{groupplot}[group style={group size=7 by 2}]
\nextgroupplot[
tick align=outside,
tick pos=left,
x grid style={darkgray176},
xmin=-1.6, xmax=22.6,
xtick style={color=black},
y grid style={darkgray176},
ymin=0, ymax=0.8587299,
ytick style={color=black}
]
\draw[draw=darkorange,fill=orange] (axis cs:-0.5,0) rectangle (axis cs:0.5,0.817838);
\draw[draw=darkorange,fill=orange] (axis cs:0.5,0) rectangle (axis cs:1.5,0.165196);
\draw[draw=darkorange,fill=orange] (axis cs:1.5,0) rectangle (axis cs:2.5,0.015945);
\draw[draw=darkorange,fill=orange] (axis cs:2.5,0) rectangle (axis cs:3.5,0.000978);
\draw[draw=darkorange,fill=orange] (axis cs:3.5,0) rectangle (axis cs:4.5,4.2e-05);
\draw[draw=darkorange,fill=orange] (axis cs:4.5,0) rectangle (axis cs:5.5,1e-06);
\draw[draw=darkorange,fill=orange] (axis cs:5.5,0) rectangle (axis cs:6.5,0);
\draw[draw=darkorange,fill=orange] (axis cs:6.5,0) rectangle (axis cs:7.5,0);
\draw[draw=darkorange,fill=orange] (axis cs:7.5,0) rectangle (axis cs:8.5,0);
\draw[draw=darkorange,fill=orange] (axis cs:8.5,0) rectangle (axis cs:9.5,0);
\draw[draw=darkorange,fill=orange] (axis cs:9.5,0) rectangle (axis cs:10.5,0);
\draw[draw=darkorange,fill=orange] (axis cs:10.5,0) rectangle (axis cs:11.5,0);
\draw[draw=darkorange,fill=orange] (axis cs:11.5,0) rectangle (axis cs:12.5,0);
\draw[draw=darkorange,fill=orange] (axis cs:12.5,0) rectangle (axis cs:13.5,0);
\draw[draw=darkorange,fill=orange] (axis cs:13.5,0) rectangle (axis cs:14.5,0);
\draw[draw=darkorange,fill=orange] (axis cs:14.5,0) rectangle (axis cs:15.5,0);
\draw[draw=darkorange,fill=orange] (axis cs:15.5,0) rectangle (axis cs:16.5,0);
\draw[draw=darkorange,fill=orange] (axis cs:16.5,0) rectangle (axis cs:17.5,0);
\draw[draw=darkorange,fill=orange] (axis cs:17.5,0) rectangle (axis cs:18.5,0);
\draw[draw=darkorange,fill=orange] (axis cs:18.5,0) rectangle (axis cs:19.5,0);
\draw[draw=darkorange,fill=orange] (axis cs:19.5,0) rectangle (axis cs:20.5,0);
\draw[draw=darkorange,fill=orange] (axis cs:20.5,0) rectangle (axis cs:21.5,0);
\draw[draw=darkslateblue10251153,fill=purple,opacity=0.5] (axis cs:-0.5,0) rectangle (axis cs:0.5,0.817785);
\draw[draw=darkslateblue10251153,fill=purple,opacity=0.5] (axis cs:0.5,0) rectangle (axis cs:1.5,0.176594);
\draw[draw=darkslateblue10251153,fill=purple,opacity=0.5] (axis cs:1.5,0) rectangle (axis cs:2.5,0.005617);
\draw[draw=darkslateblue10251153,fill=purple,opacity=0.5] (axis cs:2.5,0) rectangle (axis cs:3.5,4e-06);
\draw[draw=darkslateblue10251153,fill=purple,opacity=0.5] (axis cs:3.5,0) rectangle (axis cs:4.5,0);
\draw[draw=darkslateblue10251153,fill=purple,opacity=0.5] (axis cs:4.5,0) rectangle (axis cs:5.5,0);
\draw[draw=darkslateblue10251153,fill=purple,opacity=0.5] (axis cs:5.5,0) rectangle (axis cs:6.5,0);
\draw[draw=darkslateblue10251153,fill=purple,opacity=0.5] (axis cs:6.5,0) rectangle (axis cs:7.5,0);
\draw[draw=darkslateblue10251153,fill=purple,opacity=0.5] (axis cs:7.5,0) rectangle (axis cs:8.5,0);
\draw[draw=darkslateblue10251153,fill=purple,opacity=0.5] (axis cs:8.5,0) rectangle (axis cs:9.5,0);
\draw[draw=darkslateblue10251153,fill=purple,opacity=0.5] (axis cs:9.5,0) rectangle (axis cs:10.5,0);
\draw[draw=darkslateblue10251153,fill=purple,opacity=0.5] (axis cs:10.5,0) rectangle (axis cs:11.5,0);
\draw[draw=darkslateblue10251153,fill=purple,opacity=0.5] (axis cs:11.5,0) rectangle (axis cs:12.5,0);
\draw[draw=darkslateblue10251153,fill=purple,opacity=0.5] (axis cs:12.5,0) rectangle (axis cs:13.5,0);
\draw[draw=darkslateblue10251153,fill=purple,opacity=0.5] (axis cs:13.5,0) rectangle (axis cs:14.5,0);
\draw[draw=darkslateblue10251153,fill=purple,opacity=0.5] (axis cs:14.5,0) rectangle (axis cs:15.5,0);
\draw[draw=darkslateblue10251153,fill=purple,opacity=0.5] (axis cs:15.5,0) rectangle (axis cs:16.5,0);
\draw[draw=darkslateblue10251153,fill=purple,opacity=0.5] (axis cs:16.5,0) rectangle (axis cs:17.5,0);
\draw[draw=darkslateblue10251153,fill=purple,opacity=0.5] (axis cs:17.5,0) rectangle (axis cs:18.5,0);
\draw[draw=darkslateblue10251153,fill=purple,opacity=0.5] (axis cs:18.5,0) rectangle (axis cs:19.5,0);
\draw[draw=darkslateblue10251153,fill=purple,opacity=0.5] (axis cs:19.5,0) rectangle (axis cs:20.5,0);
\draw[draw=darkslateblue10251153,fill=purple,opacity=0.5] (axis cs:20.5,0) rectangle (axis cs:21.5,0);

\nextgroupplot[
tick align=outside,
tick pos=left,
x grid style={darkgray176},
xmin=-1.6, xmax=22.6,
xtick style={color=black},
y grid style={darkgray176},
ymin=0, ymax=0.3962742,
ytick style={color=black}
]
\draw[draw=darkorange,fill=orange] (axis cs:-0.5,0) rectangle (axis cs:0.5,0.358006);
\draw[draw=darkorange,fill=orange] (axis cs:0.5,0) rectangle (axis cs:1.5,0.377404);
\draw[draw=darkorange,fill=orange] (axis cs:1.5,0) rectangle (axis cs:2.5,0.188586);
\draw[draw=darkorange,fill=orange] (axis cs:2.5,0) rectangle (axis cs:3.5,0.060086);
\draw[draw=darkorange,fill=orange] (axis cs:3.5,0) rectangle (axis cs:4.5,0.013347);
\draw[draw=darkorange,fill=orange] (axis cs:4.5,0) rectangle (axis cs:5.5,0.002257);
\draw[draw=darkorange,fill=orange] (axis cs:5.5,0) rectangle (axis cs:6.5,0.000273);
\draw[draw=darkorange,fill=orange] (axis cs:6.5,0) rectangle (axis cs:7.5,3.8e-05);
\draw[draw=darkorange,fill=orange] (axis cs:7.5,0) rectangle (axis cs:8.5,3e-06);
\draw[draw=darkorange,fill=orange] (axis cs:8.5,0) rectangle (axis cs:9.5,0);
\draw[draw=darkorange,fill=orange] (axis cs:9.5,0) rectangle (axis cs:10.5,0);
\draw[draw=darkorange,fill=orange] (axis cs:10.5,0) rectangle (axis cs:11.5,0);
\draw[draw=darkorange,fill=orange] (axis cs:11.5,0) rectangle (axis cs:12.5,0);
\draw[draw=darkorange,fill=orange] (axis cs:12.5,0) rectangle (axis cs:13.5,0);
\draw[draw=darkorange,fill=orange] (axis cs:13.5,0) rectangle (axis cs:14.5,0);
\draw[draw=darkorange,fill=orange] (axis cs:14.5,0) rectangle (axis cs:15.5,0);
\draw[draw=darkorange,fill=orange] (axis cs:15.5,0) rectangle (axis cs:16.5,0);
\draw[draw=darkorange,fill=orange] (axis cs:16.5,0) rectangle (axis cs:17.5,0);
\draw[draw=darkorange,fill=orange] (axis cs:17.5,0) rectangle (axis cs:18.5,0);
\draw[draw=darkorange,fill=orange] (axis cs:18.5,0) rectangle (axis cs:19.5,0);
\draw[draw=darkorange,fill=orange] (axis cs:19.5,0) rectangle (axis cs:20.5,0);
\draw[draw=darkorange,fill=orange] (axis cs:20.5,0) rectangle (axis cs:21.5,0);
\draw[draw=darkslateblue10251153,fill=purple,opacity=0.5] (axis cs:-0.5,0) rectangle (axis cs:0.5,0.357675);
\draw[draw=darkslateblue10251153,fill=purple,opacity=0.5] (axis cs:0.5,0) rectangle (axis cs:1.5,0.334422);
\draw[draw=darkslateblue10251153,fill=purple,opacity=0.5] (axis cs:1.5,0) rectangle (axis cs:2.5,0.221259);
\draw[draw=darkslateblue10251153,fill=purple,opacity=0.5] (axis cs:2.5,0) rectangle (axis cs:3.5,0.076723);
\draw[draw=darkslateblue10251153,fill=purple,opacity=0.5] (axis cs:3.5,0) rectangle (axis cs:4.5,0.009499);
\draw[draw=darkslateblue10251153,fill=purple,opacity=0.5] (axis cs:4.5,0) rectangle (axis cs:5.5,0.000416);
\draw[draw=darkslateblue10251153,fill=purple,opacity=0.5] (axis cs:5.5,0) rectangle (axis cs:6.5,6e-06);
\draw[draw=darkslateblue10251153,fill=purple,opacity=0.5] (axis cs:6.5,0) rectangle (axis cs:7.5,0);
\draw[draw=darkslateblue10251153,fill=purple,opacity=0.5] (axis cs:7.5,0) rectangle (axis cs:8.5,0);
\draw[draw=darkslateblue10251153,fill=purple,opacity=0.5] (axis cs:8.5,0) rectangle (axis cs:9.5,0);
\draw[draw=darkslateblue10251153,fill=purple,opacity=0.5] (axis cs:9.5,0) rectangle (axis cs:10.5,0);
\draw[draw=darkslateblue10251153,fill=purple,opacity=0.5] (axis cs:10.5,0) rectangle (axis cs:11.5,0);
\draw[draw=darkslateblue10251153,fill=purple,opacity=0.5] (axis cs:11.5,0) rectangle (axis cs:12.5,0);
\draw[draw=darkslateblue10251153,fill=purple,opacity=0.5] (axis cs:12.5,0) rectangle (axis cs:13.5,0);
\draw[draw=darkslateblue10251153,fill=purple,opacity=0.5] (axis cs:13.5,0) rectangle (axis cs:14.5,0);
\draw[draw=darkslateblue10251153,fill=purple,opacity=0.5] (axis cs:14.5,0) rectangle (axis cs:15.5,0);
\draw[draw=darkslateblue10251153,fill=purple,opacity=0.5] (axis cs:15.5,0) rectangle (axis cs:16.5,0);
\draw[draw=darkslateblue10251153,fill=purple,opacity=0.5] (axis cs:16.5,0) rectangle (axis cs:17.5,0);
\draw[draw=darkslateblue10251153,fill=purple,opacity=0.5] (axis cs:17.5,0) rectangle (axis cs:18.5,0);
\draw[draw=darkslateblue10251153,fill=purple,opacity=0.5] (axis cs:18.5,0) rectangle (axis cs:19.5,0);
\draw[draw=darkslateblue10251153,fill=purple,opacity=0.5] (axis cs:19.5,0) rectangle (axis cs:20.5,0);
\draw[draw=darkslateblue10251153,fill=purple,opacity=0.5] (axis cs:20.5,0) rectangle (axis cs:21.5,0);

\nextgroupplot[
tick align=outside,
tick pos=left,
x grid style={darkgray176},
xmin=-1.6, xmax=22.6,
xtick style={color=black},
y grid style={darkgray176},
ymin=0, ymax=0.29913555,
ytick style={color=black}
]
\draw[draw=darkorange,fill=orange] (axis cs:-0.5,0) rectangle (axis cs:0.5,0.121628);
\draw[draw=darkorange,fill=orange] (axis cs:0.5,0) rectangle (axis cs:1.5,0.270672);
\draw[draw=darkorange,fill=orange] (axis cs:1.5,0) rectangle (axis cs:2.5,0.284891);
\draw[draw=darkorange,fill=orange] (axis cs:2.5,0) rectangle (axis cs:3.5,0.190175);
\draw[draw=darkorange,fill=orange] (axis cs:3.5,0) rectangle (axis cs:4.5,0.089547);
\draw[draw=darkorange,fill=orange] (axis cs:4.5,0) rectangle (axis cs:5.5,0.031737);
\draw[draw=darkorange,fill=orange] (axis cs:5.5,0) rectangle (axis cs:6.5,0.008964);
\draw[draw=darkorange,fill=orange] (axis cs:6.5,0) rectangle (axis cs:7.5,0.001973);
\draw[draw=darkorange,fill=orange] (axis cs:7.5,0) rectangle (axis cs:8.5,0.000351);
\draw[draw=darkorange,fill=orange] (axis cs:8.5,0) rectangle (axis cs:9.5,5.4e-05);
\draw[draw=darkorange,fill=orange] (axis cs:9.5,0) rectangle (axis cs:10.5,7e-06);
\draw[draw=darkorange,fill=orange] (axis cs:10.5,0) rectangle (axis cs:11.5,1e-06);
\draw[draw=darkorange,fill=orange] (axis cs:11.5,0) rectangle (axis cs:12.5,0);
\draw[draw=darkorange,fill=orange] (axis cs:12.5,0) rectangle (axis cs:13.5,0);
\draw[draw=darkorange,fill=orange] (axis cs:13.5,0) rectangle (axis cs:14.5,0);
\draw[draw=darkorange,fill=orange] (axis cs:14.5,0) rectangle (axis cs:15.5,0);
\draw[draw=darkorange,fill=orange] (axis cs:15.5,0) rectangle (axis cs:16.5,0);
\draw[draw=darkorange,fill=orange] (axis cs:16.5,0) rectangle (axis cs:17.5,0);
\draw[draw=darkorange,fill=orange] (axis cs:17.5,0) rectangle (axis cs:18.5,0);
\draw[draw=darkorange,fill=orange] (axis cs:18.5,0) rectangle (axis cs:19.5,0);
\draw[draw=darkorange,fill=orange] (axis cs:19.5,0) rectangle (axis cs:20.5,0);
\draw[draw=darkorange,fill=orange] (axis cs:20.5,0) rectangle (axis cs:21.5,0);
\draw[draw=darkslateblue10251153,fill=purple,opacity=0.5] (axis cs:-0.5,0) rectangle (axis cs:0.5,0.12153);
\draw[draw=darkslateblue10251153,fill=purple,opacity=0.5] (axis cs:0.5,0) rectangle (axis cs:1.5,0.216068);
\draw[draw=darkslateblue10251153,fill=purple,opacity=0.5] (axis cs:1.5,0) rectangle (axis cs:2.5,0.249803);
\draw[draw=darkslateblue10251153,fill=purple,opacity=0.5] (axis cs:2.5,0) rectangle (axis cs:3.5,0.237395);
\draw[draw=darkslateblue10251153,fill=purple,opacity=0.5] (axis cs:3.5,0) rectangle (axis cs:4.5,0.128119);
\draw[draw=darkslateblue10251153,fill=purple,opacity=0.5] (axis cs:4.5,0) rectangle (axis cs:5.5,0.039352);
\draw[draw=darkslateblue10251153,fill=purple,opacity=0.5] (axis cs:5.5,0) rectangle (axis cs:6.5,0.006957);
\draw[draw=darkslateblue10251153,fill=purple,opacity=0.5] (axis cs:6.5,0) rectangle (axis cs:7.5,0.000719);
\draw[draw=darkslateblue10251153,fill=purple,opacity=0.5] (axis cs:7.5,0) rectangle (axis cs:8.5,5.7e-05);
\draw[draw=darkslateblue10251153,fill=purple,opacity=0.5] (axis cs:8.5,0) rectangle (axis cs:9.5,0);
\draw[draw=darkslateblue10251153,fill=purple,opacity=0.5] (axis cs:9.5,0) rectangle (axis cs:10.5,0);
\draw[draw=darkslateblue10251153,fill=purple,opacity=0.5] (axis cs:10.5,0) rectangle (axis cs:11.5,0);
\draw[draw=darkslateblue10251153,fill=purple,opacity=0.5] (axis cs:11.5,0) rectangle (axis cs:12.5,0);
\draw[draw=darkslateblue10251153,fill=purple,opacity=0.5] (axis cs:12.5,0) rectangle (axis cs:13.5,0);
\draw[draw=darkslateblue10251153,fill=purple,opacity=0.5] (axis cs:13.5,0) rectangle (axis cs:14.5,0);
\draw[draw=darkslateblue10251153,fill=purple,opacity=0.5] (axis cs:14.5,0) rectangle (axis cs:15.5,0);
\draw[draw=darkslateblue10251153,fill=purple,opacity=0.5] (axis cs:15.5,0) rectangle (axis cs:16.5,0);
\draw[draw=darkslateblue10251153,fill=purple,opacity=0.5] (axis cs:16.5,0) rectangle (axis cs:17.5,0);
\draw[draw=darkslateblue10251153,fill=purple,opacity=0.5] (axis cs:17.5,0) rectangle (axis cs:18.5,0);
\draw[draw=darkslateblue10251153,fill=purple,opacity=0.5] (axis cs:18.5,0) rectangle (axis cs:19.5,0);
\draw[draw=darkslateblue10251153,fill=purple,opacity=0.5] (axis cs:19.5,0) rectangle (axis cs:20.5,0);
\draw[draw=darkslateblue10251153,fill=purple,opacity=0.5] (axis cs:20.5,0) rectangle (axis cs:21.5,0);

\nextgroupplot[
tick align=outside,
tick pos=left,
x grid style={darkgray176},
xmin=-1.6, xmax=22.6,
xtick style={color=black},
y grid style={darkgray176},
ymin=0, ymax=0.2559123,
ytick style={color=black}
]
\draw[draw=darkorange,fill=orange] (axis cs:-0.5,0) rectangle (axis cs:0.5,0.038528);
\draw[draw=darkorange,fill=orange] (axis cs:0.5,0) rectangle (axis cs:1.5,0.136709);
\draw[draw=darkorange,fill=orange] (axis cs:1.5,0) rectangle (axis cs:2.5,0.228956);
\draw[draw=darkorange,fill=orange] (axis cs:2.5,0) rectangle (axis cs:3.5,0.243726);
\draw[draw=darkorange,fill=orange] (axis cs:3.5,0) rectangle (axis cs:4.5,0.181727);
\draw[draw=darkorange,fill=orange] (axis cs:4.5,0) rectangle (axis cs:5.5,0.10295);
\draw[draw=darkorange,fill=orange] (axis cs:5.5,0) rectangle (axis cs:6.5,0.045358);
\draw[draw=darkorange,fill=orange] (axis cs:6.5,0) rectangle (axis cs:7.5,0.016204);
\draw[draw=darkorange,fill=orange] (axis cs:7.5,0) rectangle (axis cs:8.5,0.004552);
\draw[draw=darkorange,fill=orange] (axis cs:8.5,0) rectangle (axis cs:9.5,0.001064);
\draw[draw=darkorange,fill=orange] (axis cs:9.5,0) rectangle (axis cs:10.5,0.000192);
\draw[draw=darkorange,fill=orange] (axis cs:10.5,0) rectangle (axis cs:11.5,3.1e-05);
\draw[draw=darkorange,fill=orange] (axis cs:11.5,0) rectangle (axis cs:12.5,3e-06);
\draw[draw=darkorange,fill=orange] (axis cs:12.5,0) rectangle (axis cs:13.5,0);
\draw[draw=darkorange,fill=orange] (axis cs:13.5,0) rectangle (axis cs:14.5,0);
\draw[draw=darkorange,fill=orange] (axis cs:14.5,0) rectangle (axis cs:15.5,0);
\draw[draw=darkorange,fill=orange] (axis cs:15.5,0) rectangle (axis cs:16.5,0);
\draw[draw=darkorange,fill=orange] (axis cs:16.5,0) rectangle (axis cs:17.5,0);
\draw[draw=darkorange,fill=orange] (axis cs:17.5,0) rectangle (axis cs:18.5,0);
\draw[draw=darkorange,fill=orange] (axis cs:18.5,0) rectangle (axis cs:19.5,0);
\draw[draw=darkorange,fill=orange] (axis cs:19.5,0) rectangle (axis cs:20.5,0);
\draw[draw=darkorange,fill=orange] (axis cs:20.5,0) rectangle (axis cs:21.5,0);
\draw[draw=darkslateblue10251153,fill=purple,opacity=0.5] (axis cs:-0.5,0) rectangle (axis cs:0.5,0.038548);
\draw[draw=darkslateblue10251153,fill=purple,opacity=0.5] (axis cs:0.5,0) rectangle (axis cs:1.5,0.113264);
\draw[draw=darkslateblue10251153,fill=purple,opacity=0.5] (axis cs:1.5,0) rectangle (axis cs:2.5,0.168867);
\draw[draw=darkslateblue10251153,fill=purple,opacity=0.5] (axis cs:2.5,0) rectangle (axis cs:3.5,0.236719);
\draw[draw=darkslateblue10251153,fill=purple,opacity=0.5] (axis cs:3.5,0) rectangle (axis cs:4.5,0.222304);
\draw[draw=darkslateblue10251153,fill=purple,opacity=0.5] (axis cs:4.5,0) rectangle (axis cs:5.5,0.140138);
\draw[draw=darkslateblue10251153,fill=purple,opacity=0.5] (axis cs:5.5,0) rectangle (axis cs:6.5,0.059421);
\draw[draw=darkslateblue10251153,fill=purple,opacity=0.5] (axis cs:6.5,0) rectangle (axis cs:7.5,0.017083);
\draw[draw=darkslateblue10251153,fill=purple,opacity=0.5] (axis cs:7.5,0) rectangle (axis cs:8.5,0.0032);
\draw[draw=darkslateblue10251153,fill=purple,opacity=0.5] (axis cs:8.5,0) rectangle (axis cs:9.5,0.000413);
\draw[draw=darkslateblue10251153,fill=purple,opacity=0.5] (axis cs:9.5,0) rectangle (axis cs:10.5,4.2e-05);
\draw[draw=darkslateblue10251153,fill=purple,opacity=0.5] (axis cs:10.5,0) rectangle (axis cs:11.5,1e-06);
\draw[draw=darkslateblue10251153,fill=purple,opacity=0.5] (axis cs:11.5,0) rectangle (axis cs:12.5,0);
\draw[draw=darkslateblue10251153,fill=purple,opacity=0.5] (axis cs:12.5,0) rectangle (axis cs:13.5,0);
\draw[draw=darkslateblue10251153,fill=purple,opacity=0.5] (axis cs:13.5,0) rectangle (axis cs:14.5,0);
\draw[draw=darkslateblue10251153,fill=purple,opacity=0.5] (axis cs:14.5,0) rectangle (axis cs:15.5,0);
\draw[draw=darkslateblue10251153,fill=purple,opacity=0.5] (axis cs:15.5,0) rectangle (axis cs:16.5,0);
\draw[draw=darkslateblue10251153,fill=purple,opacity=0.5] (axis cs:16.5,0) rectangle (axis cs:17.5,0);
\draw[draw=darkslateblue10251153,fill=purple,opacity=0.5] (axis cs:17.5,0) rectangle (axis cs:18.5,0);
\draw[draw=darkslateblue10251153,fill=purple,opacity=0.5] (axis cs:18.5,0) rectangle (axis cs:19.5,0);
\draw[draw=darkslateblue10251153,fill=purple,opacity=0.5] (axis cs:19.5,0) rectangle (axis cs:20.5,0);
\draw[draw=darkslateblue10251153,fill=purple,opacity=0.5] (axis cs:20.5,0) rectangle (axis cs:21.5,0);

\nextgroupplot[
tick align=outside,
tick pos=left,
x grid style={darkgray176},
xmin=-1.6, xmax=22.6,
xtick style={color=black},
y grid style={darkgray176},
ymin=0, ymax=0.19059705,
ytick style={color=black}
]
\draw[draw=darkorange,fill=orange] (axis cs:-0.5,0) rectangle (axis cs:0.5,0);
\draw[draw=darkorange,fill=orange] (axis cs:0.5,0) rectangle (axis cs:1.5,2.1e-05);
\draw[draw=darkorange,fill=orange] (axis cs:1.5,0) rectangle (axis cs:2.5,0.000182);
\draw[draw=darkorange,fill=orange] (axis cs:2.5,0) rectangle (axis cs:3.5,0.001096);
\draw[draw=darkorange,fill=orange] (axis cs:3.5,0) rectangle (axis cs:4.5,0.004643);
\draw[draw=darkorange,fill=orange] (axis cs:4.5,0) rectangle (axis cs:5.5,0.014673);
\draw[draw=darkorange,fill=orange] (axis cs:5.5,0) rectangle (axis cs:6.5,0.037108);
\draw[draw=darkorange,fill=orange] (axis cs:6.5,0) rectangle (axis cs:7.5,0.073737);
\draw[draw=darkorange,fill=orange] (axis cs:7.5,0) rectangle (axis cs:8.5,0.120212);
\draw[draw=darkorange,fill=orange] (axis cs:8.5,0) rectangle (axis cs:9.5,0.159763);
\draw[draw=darkorange,fill=orange] (axis cs:9.5,0) rectangle (axis cs:10.5,0.176059);
\draw[draw=darkorange,fill=orange] (axis cs:10.5,0) rectangle (axis cs:11.5,0.160661);
\draw[draw=darkorange,fill=orange] (axis cs:11.5,0) rectangle (axis cs:12.5,0.120301);
\draw[draw=darkorange,fill=orange] (axis cs:12.5,0) rectangle (axis cs:13.5,0.074193);
\draw[draw=darkorange,fill=orange] (axis cs:13.5,0) rectangle (axis cs:14.5,0.036769);
\draw[draw=darkorange,fill=orange] (axis cs:14.5,0) rectangle (axis cs:15.5,0.014738);
\draw[draw=darkorange,fill=orange] (axis cs:15.5,0) rectangle (axis cs:16.5,0.004501);
\draw[draw=darkorange,fill=orange] (axis cs:16.5,0) rectangle (axis cs:17.5,0.001142);
\draw[draw=darkorange,fill=orange] (axis cs:17.5,0) rectangle (axis cs:18.5,0.000186);
\draw[draw=darkorange,fill=orange] (axis cs:18.5,0) rectangle (axis cs:19.5,1.3e-05);
\draw[draw=darkorange,fill=orange] (axis cs:19.5,0) rectangle (axis cs:20.5,2e-06);
\draw[draw=darkorange,fill=orange] (axis cs:20.5,0) rectangle (axis cs:21.5,0);
\draw[draw=darkslateblue10251153,fill=purple,opacity=0.5] (axis cs:-0.5,0) rectangle (axis cs:0.5,2e-06);
\draw[draw=darkslateblue10251153,fill=purple,opacity=0.5] (axis cs:0.5,0) rectangle (axis cs:1.5,4.2e-05);
\draw[draw=darkslateblue10251153,fill=purple,opacity=0.5] (axis cs:1.5,0) rectangle (axis cs:2.5,0.000265);
\draw[draw=darkslateblue10251153,fill=purple,opacity=0.5] (axis cs:2.5,0) rectangle (axis cs:3.5,0.001142);
\draw[draw=darkslateblue10251153,fill=purple,opacity=0.5] (axis cs:3.5,0) rectangle (axis cs:4.5,0.004385);
\draw[draw=darkslateblue10251153,fill=purple,opacity=0.5] (axis cs:4.5,0) rectangle (axis cs:5.5,0.013676);
\draw[draw=darkslateblue10251153,fill=purple,opacity=0.5] (axis cs:5.5,0) rectangle (axis cs:6.5,0.034388);
\draw[draw=darkslateblue10251153,fill=purple,opacity=0.5] (axis cs:6.5,0) rectangle (axis cs:7.5,0.071324);
\draw[draw=darkslateblue10251153,fill=purple,opacity=0.5] (axis cs:7.5,0) rectangle (axis cs:8.5,0.120244);
\draw[draw=darkslateblue10251153,fill=purple,opacity=0.5] (axis cs:8.5,0) rectangle (axis cs:9.5,0.164035);
\draw[draw=darkslateblue10251153,fill=purple,opacity=0.5] (axis cs:9.5,0) rectangle (axis cs:10.5,0.181521);
\draw[draw=darkslateblue10251153,fill=purple,opacity=0.5] (axis cs:10.5,0) rectangle (axis cs:11.5,0.164221);
\draw[draw=darkslateblue10251153,fill=purple,opacity=0.5] (axis cs:11.5,0) rectangle (axis cs:12.5,0.119287);
\draw[draw=darkslateblue10251153,fill=purple,opacity=0.5] (axis cs:12.5,0) rectangle (axis cs:13.5,0.071442);
\draw[draw=darkslateblue10251153,fill=purple,opacity=0.5] (axis cs:13.5,0) rectangle (axis cs:14.5,0.034324);
\draw[draw=darkslateblue10251153,fill=purple,opacity=0.5] (axis cs:14.5,0) rectangle (axis cs:15.5,0.01387);
\draw[draw=darkslateblue10251153,fill=purple,opacity=0.5] (axis cs:15.5,0) rectangle (axis cs:16.5,0.004395);
\draw[draw=darkslateblue10251153,fill=purple,opacity=0.5] (axis cs:16.5,0) rectangle (axis cs:17.5,0.001147);
\draw[draw=darkslateblue10251153,fill=purple,opacity=0.5] (axis cs:17.5,0) rectangle (axis cs:18.5,0.000253);
\draw[draw=darkslateblue10251153,fill=purple,opacity=0.5] (axis cs:18.5,0) rectangle (axis cs:19.5,2.9e-05);
\draw[draw=darkslateblue10251153,fill=purple,opacity=0.5] (axis cs:19.5,0) rectangle (axis cs:20.5,7e-06);
\draw[draw=darkslateblue10251153,fill=purple,opacity=0.5] (axis cs:20.5,0) rectangle (axis cs:21.5,1e-06);

\nextgroupplot[
tick align=outside,
tick pos=left,
x grid style={darkgray176},
xmin=-1.6, xmax=22.6,
xtick style={color=black},
y grid style={darkgray176},
ymin=0, ymax=0.30025695,
ytick style={color=black}
]
\draw[draw=darkorange,fill=orange] (axis cs:-0.5,0) rectangle (axis cs:0.5,0);
\draw[draw=darkorange,fill=orange] (axis cs:0.5,0) rectangle (axis cs:1.5,0);
\draw[draw=darkorange,fill=orange] (axis cs:1.5,0) rectangle (axis cs:2.5,0);
\draw[draw=darkorange,fill=orange] (axis cs:2.5,0) rectangle (axis cs:3.5,0);
\draw[draw=darkorange,fill=orange] (axis cs:3.5,0) rectangle (axis cs:4.5,0);
\draw[draw=darkorange,fill=orange] (axis cs:4.5,0) rectangle (axis cs:5.5,0);
\draw[draw=darkorange,fill=orange] (axis cs:5.5,0) rectangle (axis cs:6.5,0);
\draw[draw=darkorange,fill=orange] (axis cs:6.5,0) rectangle (axis cs:7.5,0);
\draw[draw=darkorange,fill=orange] (axis cs:7.5,0) rectangle (axis cs:8.5,0);
\draw[draw=darkorange,fill=orange] (axis cs:8.5,0) rectangle (axis cs:9.5,2e-06);
\draw[draw=darkorange,fill=orange] (axis cs:9.5,0) rectangle (axis cs:10.5,1.1e-05);
\draw[draw=darkorange,fill=orange] (axis cs:10.5,0) rectangle (axis cs:11.5,6.3e-05);
\draw[draw=darkorange,fill=orange] (axis cs:11.5,0) rectangle (axis cs:12.5,0.000364);
\draw[draw=darkorange,fill=orange] (axis cs:12.5,0) rectangle (axis cs:13.5,0.001924);
\draw[draw=darkorange,fill=orange] (axis cs:13.5,0) rectangle (axis cs:14.5,0.008738);
\draw[draw=darkorange,fill=orange] (axis cs:14.5,0) rectangle (axis cs:15.5,0.032122);
\draw[draw=darkorange,fill=orange] (axis cs:15.5,0) rectangle (axis cs:16.5,0.089529);
\draw[draw=darkorange,fill=orange] (axis cs:16.5,0) rectangle (axis cs:17.5,0.18991);
\draw[draw=darkorange,fill=orange] (axis cs:17.5,0) rectangle (axis cs:18.5,0.285959);
\draw[draw=darkorange,fill=orange] (axis cs:18.5,0) rectangle (axis cs:19.5,0.270063);
\draw[draw=darkorange,fill=orange] (axis cs:19.5,0) rectangle (axis cs:20.5,0.121315);
\draw[draw=darkorange,fill=orange] (axis cs:20.5,0) rectangle (axis cs:21.5,0);
\draw[draw=darkslateblue10251153,fill=purple,opacity=0.5] (axis cs:-0.5,0) rectangle (axis cs:0.5,0);
\draw[draw=darkslateblue10251153,fill=purple,opacity=0.5] (axis cs:0.5,0) rectangle (axis cs:1.5,0);
\draw[draw=darkslateblue10251153,fill=purple,opacity=0.5] (axis cs:1.5,0) rectangle (axis cs:2.5,0);
\draw[draw=darkslateblue10251153,fill=purple,opacity=0.5] (axis cs:2.5,0) rectangle (axis cs:3.5,0);
\draw[draw=darkslateblue10251153,fill=purple,opacity=0.5] (axis cs:3.5,0) rectangle (axis cs:4.5,0);
\draw[draw=darkslateblue10251153,fill=purple,opacity=0.5] (axis cs:4.5,0) rectangle (axis cs:5.5,0);
\draw[draw=darkslateblue10251153,fill=purple,opacity=0.5] (axis cs:5.5,0) rectangle (axis cs:6.5,0);
\draw[draw=darkslateblue10251153,fill=purple,opacity=0.5] (axis cs:6.5,0) rectangle (axis cs:7.5,0);
\draw[draw=darkslateblue10251153,fill=purple,opacity=0.5] (axis cs:7.5,0) rectangle (axis cs:8.5,0);
\draw[draw=darkslateblue10251153,fill=purple,opacity=0.5] (axis cs:8.5,0) rectangle (axis cs:9.5,0);
\draw[draw=darkslateblue10251153,fill=purple,opacity=0.5] (axis cs:9.5,0) rectangle (axis cs:10.5,0);
\draw[draw=darkslateblue10251153,fill=purple,opacity=0.5] (axis cs:10.5,0) rectangle (axis cs:11.5,1e-06);
\draw[draw=darkslateblue10251153,fill=purple,opacity=0.5] (axis cs:11.5,0) rectangle (axis cs:12.5,5.8e-05);
\draw[draw=darkslateblue10251153,fill=purple,opacity=0.5] (axis cs:12.5,0) rectangle (axis cs:13.5,0.000791);
\draw[draw=darkslateblue10251153,fill=purple,opacity=0.5] (axis cs:13.5,0) rectangle (axis cs:14.5,0.007931);
\draw[draw=darkslateblue10251153,fill=purple,opacity=0.5] (axis cs:14.5,0) rectangle (axis cs:15.5,0.044863);
\draw[draw=darkslateblue10251153,fill=purple,opacity=0.5] (axis cs:15.5,0) rectangle (axis cs:16.5,0.146519);
\draw[draw=darkslateblue10251153,fill=purple,opacity=0.5] (axis cs:16.5,0) rectangle (axis cs:17.5,0.268906);
\draw[draw=darkslateblue10251153,fill=purple,opacity=0.5] (axis cs:17.5,0) rectangle (axis cs:18.5,0.284243);
\draw[draw=darkslateblue10251153,fill=purple,opacity=0.5] (axis cs:18.5,0) rectangle (axis cs:19.5,0.173378);
\draw[draw=darkslateblue10251153,fill=purple,opacity=0.5] (axis cs:19.5,0) rectangle (axis cs:20.5,0.060038);
\draw[draw=darkslateblue10251153,fill=purple,opacity=0.5] (axis cs:20.5,0) rectangle (axis cs:21.5,0.013272);

\nextgroupplot[
tick align=outside,
tick pos=left,
x grid style={darkgray176},
xmin=-1.6, xmax=22.6,
xtick style={color=black},
y grid style={darkgray176},
ymin=0, ymax=0.8585556,
ytick style={color=black}
]
\draw[draw=darkorange,fill=orange] (axis cs:-0.5,0) rectangle (axis cs:0.5,0);
\draw[draw=darkorange,fill=orange] (axis cs:0.5,0) rectangle (axis cs:1.5,0);
\draw[draw=darkorange,fill=orange] (axis cs:1.5,0) rectangle (axis cs:2.5,0);
\draw[draw=darkorange,fill=orange] (axis cs:2.5,0) rectangle (axis cs:3.5,0);
\draw[draw=darkorange,fill=orange] (axis cs:3.5,0) rectangle (axis cs:4.5,0);
\draw[draw=darkorange,fill=orange] (axis cs:4.5,0) rectangle (axis cs:5.5,0);
\draw[draw=darkorange,fill=orange] (axis cs:5.5,0) rectangle (axis cs:6.5,0);
\draw[draw=darkorange,fill=orange] (axis cs:6.5,0) rectangle (axis cs:7.5,0);
\draw[draw=darkorange,fill=orange] (axis cs:7.5,0) rectangle (axis cs:8.5,0);
\draw[draw=darkorange,fill=orange] (axis cs:8.5,0) rectangle (axis cs:9.5,0);
\draw[draw=darkorange,fill=orange] (axis cs:9.5,0) rectangle (axis cs:10.5,0);
\draw[draw=darkorange,fill=orange] (axis cs:10.5,0) rectangle (axis cs:11.5,0);
\draw[draw=darkorange,fill=orange] (axis cs:11.5,0) rectangle (axis cs:12.5,0);
\draw[draw=darkorange,fill=orange] (axis cs:12.5,0) rectangle (axis cs:13.5,0);
\draw[draw=darkorange,fill=orange] (axis cs:13.5,0) rectangle (axis cs:14.5,0);
\draw[draw=darkorange,fill=orange] (axis cs:14.5,0) rectangle (axis cs:15.5,3e-06);
\draw[draw=darkorange,fill=orange] (axis cs:15.5,0) rectangle (axis cs:16.5,4.2e-05);
\draw[draw=darkorange,fill=orange] (axis cs:16.5,0) rectangle (axis cs:17.5,0.000951);
\draw[draw=darkorange,fill=orange] (axis cs:17.5,0) rectangle (axis cs:18.5,0.015929);
\draw[draw=darkorange,fill=orange] (axis cs:18.5,0) rectangle (axis cs:19.5,0.165403);
\draw[draw=darkorange,fill=orange] (axis cs:19.5,0) rectangle (axis cs:20.5,0.817672);
\draw[draw=darkorange,fill=orange] (axis cs:20.5,0) rectangle (axis cs:21.5,0);
\draw[draw=darkslateblue10251153,fill=purple,opacity=0.5] (axis cs:-0.5,0) rectangle (axis cs:0.5,0);
\draw[draw=darkslateblue10251153,fill=purple,opacity=0.5] (axis cs:0.5,0) rectangle (axis cs:1.5,0);
\draw[draw=darkslateblue10251153,fill=purple,opacity=0.5] (axis cs:1.5,0) rectangle (axis cs:2.5,0);
\draw[draw=darkslateblue10251153,fill=purple,opacity=0.5] (axis cs:2.5,0) rectangle (axis cs:3.5,0);
\draw[draw=darkslateblue10251153,fill=purple,opacity=0.5] (axis cs:3.5,0) rectangle (axis cs:4.5,0);
\draw[draw=darkslateblue10251153,fill=purple,opacity=0.5] (axis cs:4.5,0) rectangle (axis cs:5.5,0);
\draw[draw=darkslateblue10251153,fill=purple,opacity=0.5] (axis cs:5.5,0) rectangle (axis cs:6.5,0);
\draw[draw=darkslateblue10251153,fill=purple,opacity=0.5] (axis cs:6.5,0) rectangle (axis cs:7.5,0);
\draw[draw=darkslateblue10251153,fill=purple,opacity=0.5] (axis cs:7.5,0) rectangle (axis cs:8.5,0);
\draw[draw=darkslateblue10251153,fill=purple,opacity=0.5] (axis cs:8.5,0) rectangle (axis cs:9.5,0);
\draw[draw=darkslateblue10251153,fill=purple,opacity=0.5] (axis cs:9.5,0) rectangle (axis cs:10.5,0);
\draw[draw=darkslateblue10251153,fill=purple,opacity=0.5] (axis cs:10.5,0) rectangle (axis cs:11.5,0);
\draw[draw=darkslateblue10251153,fill=purple,opacity=0.5] (axis cs:11.5,0) rectangle (axis cs:12.5,0);
\draw[draw=darkslateblue10251153,fill=purple,opacity=0.5] (axis cs:12.5,0) rectangle (axis cs:13.5,0);
\draw[draw=darkslateblue10251153,fill=purple,opacity=0.5] (axis cs:13.5,0) rectangle (axis cs:14.5,0);
\draw[draw=darkslateblue10251153,fill=purple,opacity=0.5] (axis cs:14.5,0) rectangle (axis cs:15.5,0);
\draw[draw=darkslateblue10251153,fill=purple,opacity=0.5] (axis cs:15.5,0) rectangle (axis cs:16.5,0);
\draw[draw=darkslateblue10251153,fill=purple,opacity=0.5] (axis cs:16.5,0) rectangle (axis cs:17.5,1e-05);
\draw[draw=darkslateblue10251153,fill=purple,opacity=0.5] (axis cs:17.5,0) rectangle (axis cs:18.5,0.03074);
\draw[draw=darkslateblue10251153,fill=purple,opacity=0.5] (axis cs:18.5,0) rectangle (axis cs:19.5,0.637877);
\draw[draw=darkslateblue10251153,fill=purple,opacity=0.5] (axis cs:19.5,0) rectangle (axis cs:20.5,0.328358);
\draw[draw=darkslateblue10251153,fill=purple,opacity=0.5] (axis cs:20.5,0) rectangle (axis cs:21.5,0.003015);

\nextgroupplot[
tick align=outside,
tick pos=left,
x grid style={darkgray176},
xmin=-1.6, xmax=22.6,
xtick style={color=black},
y grid style={darkgray176},
ymin=0, ymax=0.8587299,
ytick style={color=black}
]
\draw[draw=darkorange,fill=orange] (axis cs:-0.5,0) rectangle (axis cs:0.5,0.817838);
\draw[draw=darkorange,fill=orange] (axis cs:0.5,0) rectangle (axis cs:1.5,0.165196);
\draw[draw=darkorange,fill=orange] (axis cs:1.5,0) rectangle (axis cs:2.5,0.015945);
\draw[draw=darkorange,fill=orange] (axis cs:2.5,0) rectangle (axis cs:3.5,0.000978);
\draw[draw=darkorange,fill=orange] (axis cs:3.5,0) rectangle (axis cs:4.5,4.2e-05);
\draw[draw=darkorange,fill=orange] (axis cs:4.5,0) rectangle (axis cs:5.5,1e-06);
\draw[draw=darkorange,fill=orange] (axis cs:5.5,0) rectangle (axis cs:6.5,0);
\draw[draw=darkorange,fill=orange] (axis cs:6.5,0) rectangle (axis cs:7.5,0);
\draw[draw=darkorange,fill=orange] (axis cs:7.5,0) rectangle (axis cs:8.5,0);
\draw[draw=darkorange,fill=orange] (axis cs:8.5,0) rectangle (axis cs:9.5,0);
\draw[draw=darkorange,fill=orange] (axis cs:9.5,0) rectangle (axis cs:10.5,0);
\draw[draw=darkorange,fill=orange] (axis cs:10.5,0) rectangle (axis cs:11.5,0);
\draw[draw=darkorange,fill=orange] (axis cs:11.5,0) rectangle (axis cs:12.5,0);
\draw[draw=darkorange,fill=orange] (axis cs:12.5,0) rectangle (axis cs:13.5,0);
\draw[draw=darkorange,fill=orange] (axis cs:13.5,0) rectangle (axis cs:14.5,0);
\draw[draw=darkorange,fill=orange] (axis cs:14.5,0) rectangle (axis cs:15.5,0);
\draw[draw=darkorange,fill=orange] (axis cs:15.5,0) rectangle (axis cs:16.5,0);
\draw[draw=darkorange,fill=orange] (axis cs:16.5,0) rectangle (axis cs:17.5,0);
\draw[draw=darkorange,fill=orange] (axis cs:17.5,0) rectangle (axis cs:18.5,0);
\draw[draw=darkorange,fill=orange] (axis cs:18.5,0) rectangle (axis cs:19.5,0);
\draw[draw=darkorange,fill=orange] (axis cs:19.5,0) rectangle (axis cs:20.5,0);
\draw[draw=darkorange,fill=orange] (axis cs:20.5,0) rectangle (axis cs:21.5,0);
\draw[draw=darkgreen,fill=green,opacity=0.5] (axis cs:-0.5,0) rectangle (axis cs:0.5,0.749269);
\draw[draw=darkgreen,fill=green,opacity=0.5] (axis cs:0.5,0) rectangle (axis cs:1.5,0.249043);
\draw[draw=darkgreen,fill=green,opacity=0.5] (axis cs:1.5,0) rectangle (axis cs:2.5,0.001688);
\draw[draw=darkgreen,fill=green,opacity=0.5] (axis cs:2.5,0) rectangle (axis cs:3.5,0);
\draw[draw=darkgreen,fill=green,opacity=0.5] (axis cs:3.5,0) rectangle (axis cs:4.5,0);
\draw[draw=darkgreen,fill=green,opacity=0.5] (axis cs:4.5,0) rectangle (axis cs:5.5,0);
\draw[draw=darkgreen,fill=green,opacity=0.5] (axis cs:5.5,0) rectangle (axis cs:6.5,0);
\draw[draw=darkgreen,fill=green,opacity=0.5] (axis cs:6.5,0) rectangle (axis cs:7.5,0);
\draw[draw=darkgreen,fill=green,opacity=0.5] (axis cs:7.5,0) rectangle (axis cs:8.5,0);
\draw[draw=darkgreen,fill=green,opacity=0.5] (axis cs:8.5,0) rectangle (axis cs:9.5,0);
\draw[draw=darkgreen,fill=green,opacity=0.5] (axis cs:9.5,0) rectangle (axis cs:10.5,0);
\draw[draw=darkgreen,fill=green,opacity=0.5] (axis cs:10.5,0) rectangle (axis cs:11.5,0);
\draw[draw=darkgreen,fill=green,opacity=0.5] (axis cs:11.5,0) rectangle (axis cs:12.5,0);
\draw[draw=darkgreen,fill=green,opacity=0.5] (axis cs:12.5,0) rectangle (axis cs:13.5,0);
\draw[draw=darkgreen,fill=green,opacity=0.5] (axis cs:13.5,0) rectangle (axis cs:14.5,0);
\draw[draw=darkgreen,fill=green,opacity=0.5] (axis cs:14.5,0) rectangle (axis cs:15.5,0);
\draw[draw=darkgreen,fill=green,opacity=0.5] (axis cs:15.5,0) rectangle (axis cs:16.5,0);
\draw[draw=darkgreen,fill=green,opacity=0.5] (axis cs:16.5,0) rectangle (axis cs:17.5,0);
\draw[draw=darkgreen,fill=green,opacity=0.5] (axis cs:17.5,0) rectangle (axis cs:18.5,0);
\draw[draw=darkgreen,fill=green,opacity=0.5] (axis cs:18.5,0) rectangle (axis cs:19.5,0);
\draw[draw=darkgreen,fill=green,opacity=0.5] (axis cs:19.5,0) rectangle (axis cs:20.5,0);
\draw[draw=darkgreen,fill=green,opacity=0.5] (axis cs:20.5,0) rectangle (axis cs:21.5,0);

\nextgroupplot[
tick align=outside,
tick pos=left,
x grid style={darkgray176},
xmin=-1.6, xmax=22.6,
xtick style={color=black},
y grid style={darkgray176},
ymin=0, ymax=0.41246205,
ytick style={color=black}
]
\draw[draw=darkorange,fill=orange] (axis cs:-0.5,0) rectangle (axis cs:0.5,0.358006);
\draw[draw=darkorange,fill=orange] (axis cs:0.5,0) rectangle (axis cs:1.5,0.377404);
\draw[draw=darkorange,fill=orange] (axis cs:1.5,0) rectangle (axis cs:2.5,0.188586);
\draw[draw=darkorange,fill=orange] (axis cs:2.5,0) rectangle (axis cs:3.5,0.060086);
\draw[draw=darkorange,fill=orange] (axis cs:3.5,0) rectangle (axis cs:4.5,0.013347);
\draw[draw=darkorange,fill=orange] (axis cs:4.5,0) rectangle (axis cs:5.5,0.002257);
\draw[draw=darkorange,fill=orange] (axis cs:5.5,0) rectangle (axis cs:6.5,0.000273);
\draw[draw=darkorange,fill=orange] (axis cs:6.5,0) rectangle (axis cs:7.5,3.8e-05);
\draw[draw=darkorange,fill=orange] (axis cs:7.5,0) rectangle (axis cs:8.5,3e-06);
\draw[draw=darkorange,fill=orange] (axis cs:8.5,0) rectangle (axis cs:9.5,0);
\draw[draw=darkorange,fill=orange] (axis cs:9.5,0) rectangle (axis cs:10.5,0);
\draw[draw=darkorange,fill=orange] (axis cs:10.5,0) rectangle (axis cs:11.5,0);
\draw[draw=darkorange,fill=orange] (axis cs:11.5,0) rectangle (axis cs:12.5,0);
\draw[draw=darkorange,fill=orange] (axis cs:12.5,0) rectangle (axis cs:13.5,0);
\draw[draw=darkorange,fill=orange] (axis cs:13.5,0) rectangle (axis cs:14.5,0);
\draw[draw=darkorange,fill=orange] (axis cs:14.5,0) rectangle (axis cs:15.5,0);
\draw[draw=darkorange,fill=orange] (axis cs:15.5,0) rectangle (axis cs:16.5,0);
\draw[draw=darkorange,fill=orange] (axis cs:16.5,0) rectangle (axis cs:17.5,0);
\draw[draw=darkorange,fill=orange] (axis cs:17.5,0) rectangle (axis cs:18.5,0);
\draw[draw=darkorange,fill=orange] (axis cs:18.5,0) rectangle (axis cs:19.5,0);
\draw[draw=darkorange,fill=orange] (axis cs:19.5,0) rectangle (axis cs:20.5,0);
\draw[draw=darkorange,fill=orange] (axis cs:20.5,0) rectangle (axis cs:21.5,0);
\draw[draw=darkgreen,fill=green,opacity=0.5] (axis cs:-0.5,0) rectangle (axis cs:0.5,0.304168);
\draw[draw=darkgreen,fill=green,opacity=0.5] (axis cs:0.5,0) rectangle (axis cs:1.5,0.392821);
\draw[draw=darkgreen,fill=green,opacity=0.5] (axis cs:1.5,0) rectangle (axis cs:2.5,0.240928);
\draw[draw=darkgreen,fill=green,opacity=0.5] (axis cs:2.5,0) rectangle (axis cs:3.5,0.056877);
\draw[draw=darkgreen,fill=green,opacity=0.5] (axis cs:3.5,0) rectangle (axis cs:4.5,0.005023);
\draw[draw=darkgreen,fill=green,opacity=0.5] (axis cs:4.5,0) rectangle (axis cs:5.5,0.000181);
\draw[draw=darkgreen,fill=green,opacity=0.5] (axis cs:5.5,0) rectangle (axis cs:6.5,2e-06);
\draw[draw=darkgreen,fill=green,opacity=0.5] (axis cs:6.5,0) rectangle (axis cs:7.5,0);
\draw[draw=darkgreen,fill=green,opacity=0.5] (axis cs:7.5,0) rectangle (axis cs:8.5,0);
\draw[draw=darkgreen,fill=green,opacity=0.5] (axis cs:8.5,0) rectangle (axis cs:9.5,0);
\draw[draw=darkgreen,fill=green,opacity=0.5] (axis cs:9.5,0) rectangle (axis cs:10.5,0);
\draw[draw=darkgreen,fill=green,opacity=0.5] (axis cs:10.5,0) rectangle (axis cs:11.5,0);
\draw[draw=darkgreen,fill=green,opacity=0.5] (axis cs:11.5,0) rectangle (axis cs:12.5,0);
\draw[draw=darkgreen,fill=green,opacity=0.5] (axis cs:12.5,0) rectangle (axis cs:13.5,0);
\draw[draw=darkgreen,fill=green,opacity=0.5] (axis cs:13.5,0) rectangle (axis cs:14.5,0);
\draw[draw=darkgreen,fill=green,opacity=0.5] (axis cs:14.5,0) rectangle (axis cs:15.5,0);
\draw[draw=darkgreen,fill=green,opacity=0.5] (axis cs:15.5,0) rectangle (axis cs:16.5,0);
\draw[draw=darkgreen,fill=green,opacity=0.5] (axis cs:16.5,0) rectangle (axis cs:17.5,0);
\draw[draw=darkgreen,fill=green,opacity=0.5] (axis cs:17.5,0) rectangle (axis cs:18.5,0);
\draw[draw=darkgreen,fill=green,opacity=0.5] (axis cs:18.5,0) rectangle (axis cs:19.5,0);
\draw[draw=darkgreen,fill=green,opacity=0.5] (axis cs:19.5,0) rectangle (axis cs:20.5,0);
\draw[draw=darkgreen,fill=green,opacity=0.5] (axis cs:20.5,0) rectangle (axis cs:21.5,0);

\nextgroupplot[
tick align=outside,
tick pos=left,
x grid style={darkgray176},
xmin=-1.6, xmax=22.6,
xtick style={color=black},
y grid style={darkgray176},
ymin=0, ymax=0.3058923,
ytick style={color=black}
]
\draw[draw=darkorange,fill=orange] (axis cs:-0.5,0) rectangle (axis cs:0.5,0.121628);
\draw[draw=darkorange,fill=orange] (axis cs:0.5,0) rectangle (axis cs:1.5,0.270672);
\draw[draw=darkorange,fill=orange] (axis cs:1.5,0) rectangle (axis cs:2.5,0.284891);
\draw[draw=darkorange,fill=orange] (axis cs:2.5,0) rectangle (axis cs:3.5,0.190175);
\draw[draw=darkorange,fill=orange] (axis cs:3.5,0) rectangle (axis cs:4.5,0.089547);
\draw[draw=darkorange,fill=orange] (axis cs:4.5,0) rectangle (axis cs:5.5,0.031737);
\draw[draw=darkorange,fill=orange] (axis cs:5.5,0) rectangle (axis cs:6.5,0.008964);
\draw[draw=darkorange,fill=orange] (axis cs:6.5,0) rectangle (axis cs:7.5,0.001973);
\draw[draw=darkorange,fill=orange] (axis cs:7.5,0) rectangle (axis cs:8.5,0.000351);
\draw[draw=darkorange,fill=orange] (axis cs:8.5,0) rectangle (axis cs:9.5,5.4e-05);
\draw[draw=darkorange,fill=orange] (axis cs:9.5,0) rectangle (axis cs:10.5,7e-06);
\draw[draw=darkorange,fill=orange] (axis cs:10.5,0) rectangle (axis cs:11.5,1e-06);
\draw[draw=darkorange,fill=orange] (axis cs:11.5,0) rectangle (axis cs:12.5,0);
\draw[draw=darkorange,fill=orange] (axis cs:12.5,0) rectangle (axis cs:13.5,0);
\draw[draw=darkorange,fill=orange] (axis cs:13.5,0) rectangle (axis cs:14.5,0);
\draw[draw=darkorange,fill=orange] (axis cs:14.5,0) rectangle (axis cs:15.5,0);
\draw[draw=darkorange,fill=orange] (axis cs:15.5,0) rectangle (axis cs:16.5,0);
\draw[draw=darkorange,fill=orange] (axis cs:16.5,0) rectangle (axis cs:17.5,0);
\draw[draw=darkorange,fill=orange] (axis cs:17.5,0) rectangle (axis cs:18.5,0);
\draw[draw=darkorange,fill=orange] (axis cs:18.5,0) rectangle (axis cs:19.5,0);
\draw[draw=darkorange,fill=orange] (axis cs:19.5,0) rectangle (axis cs:20.5,0);
\draw[draw=darkorange,fill=orange] (axis cs:20.5,0) rectangle (axis cs:21.5,0);
\draw[draw=darkgreen,fill=green,opacity=0.5] (axis cs:-0.5,0) rectangle (axis cs:0.5,0.131937);
\draw[draw=darkgreen,fill=green,opacity=0.5] (axis cs:0.5,0) rectangle (axis cs:1.5,0.22208);
\draw[draw=darkgreen,fill=green,opacity=0.5] (axis cs:1.5,0) rectangle (axis cs:2.5,0.291326);
\draw[draw=darkgreen,fill=green,opacity=0.5] (axis cs:2.5,0) rectangle (axis cs:3.5,0.223201);
\draw[draw=darkgreen,fill=green,opacity=0.5] (axis cs:3.5,0) rectangle (axis cs:4.5,0.100304);
\draw[draw=darkgreen,fill=green,opacity=0.5] (axis cs:4.5,0) rectangle (axis cs:5.5,0.026709);
\draw[draw=darkgreen,fill=green,opacity=0.5] (axis cs:5.5,0) rectangle (axis cs:6.5,0.00403);
\draw[draw=darkgreen,fill=green,opacity=0.5] (axis cs:6.5,0) rectangle (axis cs:7.5,0.000398);
\draw[draw=darkgreen,fill=green,opacity=0.5] (axis cs:7.5,0) rectangle (axis cs:8.5,1.5e-05);
\draw[draw=darkgreen,fill=green,opacity=0.5] (axis cs:8.5,0) rectangle (axis cs:9.5,0);
\draw[draw=darkgreen,fill=green,opacity=0.5] (axis cs:9.5,0) rectangle (axis cs:10.5,0);
\draw[draw=darkgreen,fill=green,opacity=0.5] (axis cs:10.5,0) rectangle (axis cs:11.5,0);
\draw[draw=darkgreen,fill=green,opacity=0.5] (axis cs:11.5,0) rectangle (axis cs:12.5,0);
\draw[draw=darkgreen,fill=green,opacity=0.5] (axis cs:12.5,0) rectangle (axis cs:13.5,0);
\draw[draw=darkgreen,fill=green,opacity=0.5] (axis cs:13.5,0) rectangle (axis cs:14.5,0);
\draw[draw=darkgreen,fill=green,opacity=0.5] (axis cs:14.5,0) rectangle (axis cs:15.5,0);
\draw[draw=darkgreen,fill=green,opacity=0.5] (axis cs:15.5,0) rectangle (axis cs:16.5,0);
\draw[draw=darkgreen,fill=green,opacity=0.5] (axis cs:16.5,0) rectangle (axis cs:17.5,0);
\draw[draw=darkgreen,fill=green,opacity=0.5] (axis cs:17.5,0) rectangle (axis cs:18.5,0);
\draw[draw=darkgreen,fill=green,opacity=0.5] (axis cs:18.5,0) rectangle (axis cs:19.5,0);
\draw[draw=darkgreen,fill=green,opacity=0.5] (axis cs:19.5,0) rectangle (axis cs:20.5,0);
\draw[draw=darkgreen,fill=green,opacity=0.5] (axis cs:20.5,0) rectangle (axis cs:21.5,0);

\nextgroupplot[
tick align=outside,
tick pos=left,
x grid style={darkgray176},
xmin=-1.6, xmax=22.6,
xtick style={color=black},
y grid style={darkgray176},
ymin=0, ymax=0.2580984,
ytick style={color=black}
]
\draw[draw=darkorange,fill=orange] (axis cs:-0.5,0) rectangle (axis cs:0.5,0.038528);
\draw[draw=darkorange,fill=orange] (axis cs:0.5,0) rectangle (axis cs:1.5,0.136709);
\draw[draw=darkorange,fill=orange] (axis cs:1.5,0) rectangle (axis cs:2.5,0.228956);
\draw[draw=darkorange,fill=orange] (axis cs:2.5,0) rectangle (axis cs:3.5,0.243726);
\draw[draw=darkorange,fill=orange] (axis cs:3.5,0) rectangle (axis cs:4.5,0.181727);
\draw[draw=darkorange,fill=orange] (axis cs:4.5,0) rectangle (axis cs:5.5,0.10295);
\draw[draw=darkorange,fill=orange] (axis cs:5.5,0) rectangle (axis cs:6.5,0.045358);
\draw[draw=darkorange,fill=orange] (axis cs:6.5,0) rectangle (axis cs:7.5,0.016204);
\draw[draw=darkorange,fill=orange] (axis cs:7.5,0) rectangle (axis cs:8.5,0.004552);
\draw[draw=darkorange,fill=orange] (axis cs:8.5,0) rectangle (axis cs:9.5,0.001064);
\draw[draw=darkorange,fill=orange] (axis cs:9.5,0) rectangle (axis cs:10.5,0.000192);
\draw[draw=darkorange,fill=orange] (axis cs:10.5,0) rectangle (axis cs:11.5,3.1e-05);
\draw[draw=darkorange,fill=orange] (axis cs:11.5,0) rectangle (axis cs:12.5,3e-06);
\draw[draw=darkorange,fill=orange] (axis cs:12.5,0) rectangle (axis cs:13.5,0);
\draw[draw=darkorange,fill=orange] (axis cs:13.5,0) rectangle (axis cs:14.5,0);
\draw[draw=darkorange,fill=orange] (axis cs:14.5,0) rectangle (axis cs:15.5,0);
\draw[draw=darkorange,fill=orange] (axis cs:15.5,0) rectangle (axis cs:16.5,0);
\draw[draw=darkorange,fill=orange] (axis cs:16.5,0) rectangle (axis cs:17.5,0);
\draw[draw=darkorange,fill=orange] (axis cs:17.5,0) rectangle (axis cs:18.5,0);
\draw[draw=darkorange,fill=orange] (axis cs:18.5,0) rectangle (axis cs:19.5,0);
\draw[draw=darkorange,fill=orange] (axis cs:19.5,0) rectangle (axis cs:20.5,0);
\draw[draw=darkorange,fill=orange] (axis cs:20.5,0) rectangle (axis cs:21.5,0);
\draw[draw=darkgreen,fill=green,opacity=0.5] (axis cs:-0.5,0) rectangle (axis cs:0.5,0.058463);
\draw[draw=darkgreen,fill=green,opacity=0.5] (axis cs:0.5,0) rectangle (axis cs:1.5,0.114754);
\draw[draw=darkgreen,fill=green,opacity=0.5] (axis cs:1.5,0) rectangle (axis cs:2.5,0.203796);
\draw[draw=darkgreen,fill=green,opacity=0.5] (axis cs:2.5,0) rectangle (axis cs:3.5,0.245808);
\draw[draw=darkgreen,fill=green,opacity=0.5] (axis cs:3.5,0) rectangle (axis cs:4.5,0.203517);
\draw[draw=darkgreen,fill=green,opacity=0.5] (axis cs:4.5,0) rectangle (axis cs:5.5,0.11509);
\draw[draw=darkgreen,fill=green,opacity=0.5] (axis cs:5.5,0) rectangle (axis cs:6.5,0.044316);
\draw[draw=darkgreen,fill=green,opacity=0.5] (axis cs:6.5,0) rectangle (axis cs:7.5,0.011835);
\draw[draw=darkgreen,fill=green,opacity=0.5] (axis cs:7.5,0) rectangle (axis cs:8.5,0.002131);
\draw[draw=darkgreen,fill=green,opacity=0.5] (axis cs:8.5,0) rectangle (axis cs:9.5,0.000268);
\draw[draw=darkgreen,fill=green,opacity=0.5] (axis cs:9.5,0) rectangle (axis cs:10.5,2e-05);
\draw[draw=darkgreen,fill=green,opacity=0.5] (axis cs:10.5,0) rectangle (axis cs:11.5,2e-06);
\draw[draw=darkgreen,fill=green,opacity=0.5] (axis cs:11.5,0) rectangle (axis cs:12.5,0);
\draw[draw=darkgreen,fill=green,opacity=0.5] (axis cs:12.5,0) rectangle (axis cs:13.5,0);
\draw[draw=darkgreen,fill=green,opacity=0.5] (axis cs:13.5,0) rectangle (axis cs:14.5,0);
\draw[draw=darkgreen,fill=green,opacity=0.5] (axis cs:14.5,0) rectangle (axis cs:15.5,0);
\draw[draw=darkgreen,fill=green,opacity=0.5] (axis cs:15.5,0) rectangle (axis cs:16.5,0);
\draw[draw=darkgreen,fill=green,opacity=0.5] (axis cs:16.5,0) rectangle (axis cs:17.5,0);
\draw[draw=darkgreen,fill=green,opacity=0.5] (axis cs:17.5,0) rectangle (axis cs:18.5,0);
\draw[draw=darkgreen,fill=green,opacity=0.5] (axis cs:18.5,0) rectangle (axis cs:19.5,0);
\draw[draw=darkgreen,fill=green,opacity=0.5] (axis cs:19.5,0) rectangle (axis cs:20.5,0);
\draw[draw=darkgreen,fill=green,opacity=0.5] (axis cs:20.5,0) rectangle (axis cs:21.5,0);

\nextgroupplot[
tick align=outside,
tick pos=left,
x grid style={darkgray176},
xmin=-1.6, xmax=22.6,
xtick style={color=black},
y grid style={darkgray176},
ymin=0, ymax=0.1852557,
ytick style={color=black}
]
\draw[draw=darkorange,fill=orange] (axis cs:-0.5,0) rectangle (axis cs:0.5,0);
\draw[draw=darkorange,fill=orange] (axis cs:0.5,0) rectangle (axis cs:1.5,2.1e-05);
\draw[draw=darkorange,fill=orange] (axis cs:1.5,0) rectangle (axis cs:2.5,0.000182);
\draw[draw=darkorange,fill=orange] (axis cs:2.5,0) rectangle (axis cs:3.5,0.001096);
\draw[draw=darkorange,fill=orange] (axis cs:3.5,0) rectangle (axis cs:4.5,0.004643);
\draw[draw=darkorange,fill=orange] (axis cs:4.5,0) rectangle (axis cs:5.5,0.014673);
\draw[draw=darkorange,fill=orange] (axis cs:5.5,0) rectangle (axis cs:6.5,0.037108);
\draw[draw=darkorange,fill=orange] (axis cs:6.5,0) rectangle (axis cs:7.5,0.073737);
\draw[draw=darkorange,fill=orange] (axis cs:7.5,0) rectangle (axis cs:8.5,0.120212);
\draw[draw=darkorange,fill=orange] (axis cs:8.5,0) rectangle (axis cs:9.5,0.159763);
\draw[draw=darkorange,fill=orange] (axis cs:9.5,0) rectangle (axis cs:10.5,0.176059);
\draw[draw=darkorange,fill=orange] (axis cs:10.5,0) rectangle (axis cs:11.5,0.160661);
\draw[draw=darkorange,fill=orange] (axis cs:11.5,0) rectangle (axis cs:12.5,0.120301);
\draw[draw=darkorange,fill=orange] (axis cs:12.5,0) rectangle (axis cs:13.5,0.074193);
\draw[draw=darkorange,fill=orange] (axis cs:13.5,0) rectangle (axis cs:14.5,0.036769);
\draw[draw=darkorange,fill=orange] (axis cs:14.5,0) rectangle (axis cs:15.5,0.014738);
\draw[draw=darkorange,fill=orange] (axis cs:15.5,0) rectangle (axis cs:16.5,0.004501);
\draw[draw=darkorange,fill=orange] (axis cs:16.5,0) rectangle (axis cs:17.5,0.001142);
\draw[draw=darkorange,fill=orange] (axis cs:17.5,0) rectangle (axis cs:18.5,0.000186);
\draw[draw=darkorange,fill=orange] (axis cs:18.5,0) rectangle (axis cs:19.5,1.3e-05);
\draw[draw=darkorange,fill=orange] (axis cs:19.5,0) rectangle (axis cs:20.5,2e-06);
\draw[draw=darkorange,fill=orange] (axis cs:20.5,0) rectangle (axis cs:21.5,0);
\draw[draw=darkgreen,fill=green,opacity=0.5] (axis cs:-0.5,0) rectangle (axis cs:0.5,1.3e-05);
\draw[draw=darkgreen,fill=green,opacity=0.5] (axis cs:0.5,0) rectangle (axis cs:1.5,6.2e-05);
\draw[draw=darkgreen,fill=green,opacity=0.5] (axis cs:1.5,0) rectangle (axis cs:2.5,0.000318);
\draw[draw=darkgreen,fill=green,opacity=0.5] (axis cs:2.5,0) rectangle (axis cs:3.5,0.001462);
\draw[draw=darkgreen,fill=green,opacity=0.5] (axis cs:3.5,0) rectangle (axis cs:4.5,0.00511);
\draw[draw=darkgreen,fill=green,opacity=0.5] (axis cs:4.5,0) rectangle (axis cs:5.5,0.015121);
\draw[draw=darkgreen,fill=green,opacity=0.5] (axis cs:5.5,0) rectangle (axis cs:6.5,0.03694);
\draw[draw=darkgreen,fill=green,opacity=0.5] (axis cs:6.5,0) rectangle (axis cs:7.5,0.072292);
\draw[draw=darkgreen,fill=green,opacity=0.5] (axis cs:7.5,0) rectangle (axis cs:8.5,0.120065);
\draw[draw=darkgreen,fill=green,opacity=0.5] (axis cs:8.5,0) rectangle (axis cs:9.5,0.160496);
\draw[draw=darkgreen,fill=green,opacity=0.5] (axis cs:9.5,0) rectangle (axis cs:10.5,0.176434);
\draw[draw=darkgreen,fill=green,opacity=0.5] (axis cs:10.5,0) rectangle (axis cs:11.5,0.160647);
\draw[draw=darkgreen,fill=green,opacity=0.5] (axis cs:11.5,0) rectangle (axis cs:12.5,0.118966);
\draw[draw=darkgreen,fill=green,opacity=0.5] (axis cs:12.5,0) rectangle (axis cs:13.5,0.073177);
\draw[draw=darkgreen,fill=green,opacity=0.5] (axis cs:13.5,0) rectangle (axis cs:14.5,0.036848);
\draw[draw=darkgreen,fill=green,opacity=0.5] (axis cs:14.5,0) rectangle (axis cs:15.5,0.015038);
\draw[draw=darkgreen,fill=green,opacity=0.5] (axis cs:15.5,0) rectangle (axis cs:16.5,0.005157);
\draw[draw=darkgreen,fill=green,opacity=0.5] (axis cs:16.5,0) rectangle (axis cs:17.5,0.001427);
\draw[draw=darkgreen,fill=green,opacity=0.5] (axis cs:17.5,0) rectangle (axis cs:18.5,0.000355);
\draw[draw=darkgreen,fill=green,opacity=0.5] (axis cs:18.5,0) rectangle (axis cs:19.5,6.7e-05);
\draw[draw=darkgreen,fill=green,opacity=0.5] (axis cs:19.5,0) rectangle (axis cs:20.5,5e-06);
\draw[draw=darkgreen,fill=green,opacity=0.5] (axis cs:20.5,0) rectangle (axis cs:21.5,0);

\nextgroupplot[
tick align=outside,
tick pos=left,
x grid style={darkgray176},
xmin=-1.6, xmax=22.6,
xtick style={color=black},
y grid style={darkgray176},
ymin=0, ymax=0.3056235,
ytick style={color=black}
]
\draw[draw=darkorange,fill=orange] (axis cs:-0.5,0) rectangle (axis cs:0.5,0);
\draw[draw=darkorange,fill=orange] (axis cs:0.5,0) rectangle (axis cs:1.5,0);
\draw[draw=darkorange,fill=orange] (axis cs:1.5,0) rectangle (axis cs:2.5,0);
\draw[draw=darkorange,fill=orange] (axis cs:2.5,0) rectangle (axis cs:3.5,0);
\draw[draw=darkorange,fill=orange] (axis cs:3.5,0) rectangle (axis cs:4.5,0);
\draw[draw=darkorange,fill=orange] (axis cs:4.5,0) rectangle (axis cs:5.5,0);
\draw[draw=darkorange,fill=orange] (axis cs:5.5,0) rectangle (axis cs:6.5,0);
\draw[draw=darkorange,fill=orange] (axis cs:6.5,0) rectangle (axis cs:7.5,0);
\draw[draw=darkorange,fill=orange] (axis cs:7.5,0) rectangle (axis cs:8.5,0);
\draw[draw=darkorange,fill=orange] (axis cs:8.5,0) rectangle (axis cs:9.5,2e-06);
\draw[draw=darkorange,fill=orange] (axis cs:9.5,0) rectangle (axis cs:10.5,1.1e-05);
\draw[draw=darkorange,fill=orange] (axis cs:10.5,0) rectangle (axis cs:11.5,6.3e-05);
\draw[draw=darkorange,fill=orange] (axis cs:11.5,0) rectangle (axis cs:12.5,0.000364);
\draw[draw=darkorange,fill=orange] (axis cs:12.5,0) rectangle (axis cs:13.5,0.001924);
\draw[draw=darkorange,fill=orange] (axis cs:13.5,0) rectangle (axis cs:14.5,0.008738);
\draw[draw=darkorange,fill=orange] (axis cs:14.5,0) rectangle (axis cs:15.5,0.032122);
\draw[draw=darkorange,fill=orange] (axis cs:15.5,0) rectangle (axis cs:16.5,0.089529);
\draw[draw=darkorange,fill=orange] (axis cs:16.5,0) rectangle (axis cs:17.5,0.18991);
\draw[draw=darkorange,fill=orange] (axis cs:17.5,0) rectangle (axis cs:18.5,0.285959);
\draw[draw=darkorange,fill=orange] (axis cs:18.5,0) rectangle (axis cs:19.5,0.270063);
\draw[draw=darkorange,fill=orange] (axis cs:19.5,0) rectangle (axis cs:20.5,0.121315);
\draw[draw=darkorange,fill=orange] (axis cs:20.5,0) rectangle (axis cs:21.5,0);
\draw[draw=darkgreen,fill=green,opacity=0.5] (axis cs:-0.5,0) rectangle (axis cs:0.5,0);
\draw[draw=darkgreen,fill=green,opacity=0.5] (axis cs:0.5,0) rectangle (axis cs:1.5,0);
\draw[draw=darkgreen,fill=green,opacity=0.5] (axis cs:1.5,0) rectangle (axis cs:2.5,0);
\draw[draw=darkgreen,fill=green,opacity=0.5] (axis cs:2.5,0) rectangle (axis cs:3.5,0);
\draw[draw=darkgreen,fill=green,opacity=0.5] (axis cs:3.5,0) rectangle (axis cs:4.5,0);
\draw[draw=darkgreen,fill=green,opacity=0.5] (axis cs:4.5,0) rectangle (axis cs:5.5,0);
\draw[draw=darkgreen,fill=green,opacity=0.5] (axis cs:5.5,0) rectangle (axis cs:6.5,0);
\draw[draw=darkgreen,fill=green,opacity=0.5] (axis cs:6.5,0) rectangle (axis cs:7.5,0);
\draw[draw=darkgreen,fill=green,opacity=0.5] (axis cs:7.5,0) rectangle (axis cs:8.5,0);
\draw[draw=darkgreen,fill=green,opacity=0.5] (axis cs:8.5,0) rectangle (axis cs:9.5,0);
\draw[draw=darkgreen,fill=green,opacity=0.5] (axis cs:9.5,0) rectangle (axis cs:10.5,0);
\draw[draw=darkgreen,fill=green,opacity=0.5] (axis cs:10.5,0) rectangle (axis cs:11.5,1e-06);
\draw[draw=darkgreen,fill=green,opacity=0.5] (axis cs:11.5,0) rectangle (axis cs:12.5,2e-05);
\draw[draw=darkgreen,fill=green,opacity=0.5] (axis cs:12.5,0) rectangle (axis cs:13.5,0.000362);
\draw[draw=darkgreen,fill=green,opacity=0.5] (axis cs:13.5,0) rectangle (axis cs:14.5,0.004103);
\draw[draw=darkgreen,fill=green,opacity=0.5] (axis cs:14.5,0) rectangle (axis cs:15.5,0.026226);
\draw[draw=darkgreen,fill=green,opacity=0.5] (axis cs:15.5,0) rectangle (axis cs:16.5,0.100605);
\draw[draw=darkgreen,fill=green,opacity=0.5] (axis cs:16.5,0) rectangle (axis cs:17.5,0.222503);
\draw[draw=darkgreen,fill=green,opacity=0.5] (axis cs:17.5,0) rectangle (axis cs:18.5,0.29107);
\draw[draw=darkgreen,fill=green,opacity=0.5] (axis cs:18.5,0) rectangle (axis cs:19.5,0.223284);
\draw[draw=darkgreen,fill=green,opacity=0.5] (axis cs:19.5,0) rectangle (axis cs:20.5,0.131826);
\draw[draw=darkgreen,fill=green,opacity=0.5] (axis cs:20.5,0) rectangle (axis cs:21.5,0);

\nextgroupplot[
tick align=outside,
tick pos=left,
x grid style={darkgray176},
xmin=-1.6, xmax=22.6,
xtick style={color=black},
y grid style={darkgray176},
ymin=0, ymax=0.8585556,
ytick style={color=black}
]
\draw[draw=darkorange,fill=orange] (axis cs:-0.5,0) rectangle (axis cs:0.5,0);
\draw[draw=darkorange,fill=orange] (axis cs:0.5,0) rectangle (axis cs:1.5,0);
\draw[draw=darkorange,fill=orange] (axis cs:1.5,0) rectangle (axis cs:2.5,0);
\draw[draw=darkorange,fill=orange] (axis cs:2.5,0) rectangle (axis cs:3.5,0);
\draw[draw=darkorange,fill=orange] (axis cs:3.5,0) rectangle (axis cs:4.5,0);
\draw[draw=darkorange,fill=orange] (axis cs:4.5,0) rectangle (axis cs:5.5,0);
\draw[draw=darkorange,fill=orange] (axis cs:5.5,0) rectangle (axis cs:6.5,0);
\draw[draw=darkorange,fill=orange] (axis cs:6.5,0) rectangle (axis cs:7.5,0);
\draw[draw=darkorange,fill=orange] (axis cs:7.5,0) rectangle (axis cs:8.5,0);
\draw[draw=darkorange,fill=orange] (axis cs:8.5,0) rectangle (axis cs:9.5,0);
\draw[draw=darkorange,fill=orange] (axis cs:9.5,0) rectangle (axis cs:10.5,0);
\draw[draw=darkorange,fill=orange] (axis cs:10.5,0) rectangle (axis cs:11.5,0);
\draw[draw=darkorange,fill=orange] (axis cs:11.5,0) rectangle (axis cs:12.5,0);
\draw[draw=darkorange,fill=orange] (axis cs:12.5,0) rectangle (axis cs:13.5,0);
\draw[draw=darkorange,fill=orange] (axis cs:13.5,0) rectangle (axis cs:14.5,0);
\draw[draw=darkorange,fill=orange] (axis cs:14.5,0) rectangle (axis cs:15.5,3e-06);
\draw[draw=darkorange,fill=orange] (axis cs:15.5,0) rectangle (axis cs:16.5,4.2e-05);
\draw[draw=darkorange,fill=orange] (axis cs:16.5,0) rectangle (axis cs:17.5,0.000951);
\draw[draw=darkorange,fill=orange] (axis cs:17.5,0) rectangle (axis cs:18.5,0.015929);
\draw[draw=darkorange,fill=orange] (axis cs:18.5,0) rectangle (axis cs:19.5,0.165403);
\draw[draw=darkorange,fill=orange] (axis cs:19.5,0) rectangle (axis cs:20.5,0.817672);
\draw[draw=darkorange,fill=orange] (axis cs:20.5,0) rectangle (axis cs:21.5,0);
\draw[draw=darkgreen,fill=green,opacity=0.5] (axis cs:-0.5,0) rectangle (axis cs:0.5,0);
\draw[draw=darkgreen,fill=green,opacity=0.5] (axis cs:0.5,0) rectangle (axis cs:1.5,0);
\draw[draw=darkgreen,fill=green,opacity=0.5] (axis cs:1.5,0) rectangle (axis cs:2.5,0);
\draw[draw=darkgreen,fill=green,opacity=0.5] (axis cs:2.5,0) rectangle (axis cs:3.5,0);
\draw[draw=darkgreen,fill=green,opacity=0.5] (axis cs:3.5,0) rectangle (axis cs:4.5,0);
\draw[draw=darkgreen,fill=green,opacity=0.5] (axis cs:4.5,0) rectangle (axis cs:5.5,0);
\draw[draw=darkgreen,fill=green,opacity=0.5] (axis cs:5.5,0) rectangle (axis cs:6.5,0);
\draw[draw=darkgreen,fill=green,opacity=0.5] (axis cs:6.5,0) rectangle (axis cs:7.5,0);
\draw[draw=darkgreen,fill=green,opacity=0.5] (axis cs:7.5,0) rectangle (axis cs:8.5,0);
\draw[draw=darkgreen,fill=green,opacity=0.5] (axis cs:8.5,0) rectangle (axis cs:9.5,0);
\draw[draw=darkgreen,fill=green,opacity=0.5] (axis cs:9.5,0) rectangle (axis cs:10.5,0);
\draw[draw=darkgreen,fill=green,opacity=0.5] (axis cs:10.5,0) rectangle (axis cs:11.5,0);
\draw[draw=darkgreen,fill=green,opacity=0.5] (axis cs:11.5,0) rectangle (axis cs:12.5,0);
\draw[draw=darkgreen,fill=green,opacity=0.5] (axis cs:12.5,0) rectangle (axis cs:13.5,0);
\draw[draw=darkgreen,fill=green,opacity=0.5] (axis cs:13.5,0) rectangle (axis cs:14.5,0);
\draw[draw=darkgreen,fill=green,opacity=0.5] (axis cs:14.5,0) rectangle (axis cs:15.5,0);
\draw[draw=darkgreen,fill=green,opacity=0.5] (axis cs:15.5,0) rectangle (axis cs:16.5,0);
\draw[draw=darkgreen,fill=green,opacity=0.5] (axis cs:16.5,0) rectangle (axis cs:17.5,0);
\draw[draw=darkgreen,fill=green,opacity=0.5] (axis cs:17.5,0) rectangle (axis cs:18.5,0.001804);
\draw[draw=darkgreen,fill=green,opacity=0.5] (axis cs:18.5,0) rectangle (axis cs:19.5,0.248249);
\draw[draw=darkgreen,fill=green,opacity=0.5] (axis cs:19.5,0) rectangle (axis cs:20.5,0.749947);
\draw[draw=darkgreen,fill=green,opacity=0.5] (axis cs:20.5,0) rectangle (axis cs:21.5,0);
\end{groupplot}

\end{tikzpicture}
}
    \caption{Comparison of the binomial approximation by \textcite{deliot2023} (top) and a clipped normal approximation with $\mu=N\,p$ and $\sigma^2=N\,p\,(1-p)$ (bottom) to the respective binomial distribution (orange). Every distribution was sampled 1000000 times with $N=20$ and $p \in \left\{0.01, 0.05, 0.1, 0.15, 0.5, 0.9, 0.99\right\}$ from left to right. For $p=0.9$ and $p=0.99$ overshooting behavior can be observed at $X=21$. \label{fig:approximation}}
\end{figure*}

As elaborated in \cref{sec:discretization} our algorithm requires $4\times3\times4$ samples of $4\times3$ different parametrizations of the binomial distribution per pixel, therefore an efficient method for drawing these samples is required.
%
Unfortunately, common strategies like inversion sampling can not easily be applied to the binomial distribution due to its discrete nature.

\paragraph{Classical Binomial Approximation}
%
\textcite{zirr2016} therefore approximate the binomial distribution for large expectation $\expectation{b(n,p)}=np>5$ with the normal distribution $\mathcal{N}(\mu,\sigma)$ with equal mean $\mu=\expectation{b(n,p)}=np$ and variance $\sigma^2=\variance{b(n,p)}=np(1-p)$.
%
Then for the normal distribution $\mathcal{N}(\mu,\sigma)$ efficient approximations of the inverse cumulative distribution function \parencite[e.g.,][]{giles2012} and even direct sampling methods \parencite[e.g.,][]{box1958} are known.
%
For $np \le 5$ however, they rely on manually evaluating the binomial distribution which is rather expensive.


\paragraph{Gated Gaussian Binomial Approximation}
%
To achieve stable performance and to obviate extensive branching \textcite{deliot2023} present a new approach that is likewise based on the normal approximation but splits the evaluation of the binomial distribution into two independent random trials:
%
\begin{equation}
    b(N,p) \approx H(\theta_0 - \probability{X>0}) \cdot \lfloor 1 + \mathcal{N}_{\theta_1}(\mu, \sigma)\rfloor,
\end{equation}
%
thereby $H$ is the Heaviside step function, $\theta_0, \theta_1 \in \interval{0}{1}$ are uniform random variables and $\mathcal{N}_{\theta_1}(\mu, \sigma) \in \interval{0}{N}$ is a clipped normally distributed random variable, which can again be generated by inversion sampling \parencite[e.g.,][]{giles2012} or directly \parencite[e.g.,][]{box1958}.
%
Therefore, $H(\theta_0 - \probability{X>0})$ with $\probability{X>0}=1-\probability{X=0}=1-(1-p)^N$ is a random trial testing for at least one success in the evaluation of $b(N,p)$, and if this test is successful we approximate $b(N,p)$ by $1+b(N-1,p) \approx \lfloor 1 + \mathcal{N}_{\theta_1}(\mu, \sigma) \rfloor$ with $\mu=(N-1) \cdot p$ and $\sigma^2=(N-1) \cdot p \cdot (1-p)$.
%
% Let $X$ be the total number of successes, $Y=1$, if we have at least one success, $Y=0$ otherwise, and $Z$ the number of successes if we ignore the first successful trial. Then
% %
% \begin{equation}
%     X = Y \cdot (1 + Z)
% \end{equation}

This however remains only an approximation as shown in \cref{fig:approximation} and can even overshoot the number of successes to $N+1$, but as it ensures correct probability $\probability{X=0}$ for getting zero successes it performs significantly better than a normal distribution approximation for very small probability of success $p$, which commonly occurs in our shading model for low roughness.
% TODO: Maybe add figure
%
The overshooting behavior may even be desirable, as it simulates the glints sometimes being slightly shinier than the smooth base model and thus the material appears glintier (see \cref{fig:overshooting}).
%
The main advantage over the approximation by \textcite{zirr2016} however is the avoidance of loops and branches which makes this method easily parallelizable and results in stable performance.

Additionally, \textcite{deliot2023} modify the gating in their implementation from a hard Heaviside function to a slightly sloped step function to make the glints' appearance smoother.

To efficiently draw uniform random numbers and seeds on the GPU we use the recommendations by \textcite{jarzynski2020} in our implementation.